%-*-tex-*-
% Copyright Michael J. Ferguson, INRS-Telecommunications
% All rights reserved. 
% This file now uses cm fonts where available. 


% ======= INRS Fonts ========
% This file adds the fonts used in INRSTeX. It also adds the spacing 
% related to these fonts. The important section is the definitions of
% the font families. The intent is to produce a font family with all
% related fonts. In the same manner as the TeXbook, \tenpoint indicates
% the 10pt family and \rm would be a Roman font in that family. The
% additions here are the special fonts for section heads, footnotes, ...
% The section head fonts, for example \sheadfont is changed only if 
% the font family is called within the command \documentstyle. This 
% means that it is possible to change the font internal to a 
% document without disturbing the appearance of the section heads or
% page numbering. 

% In addition to this file, there is an spfont.tex. This file allows for 
% the scaling of a font family. The basic family is \tenpoint. It can be 
% scaled to the extent that there exist scaled printable forms of the font.

% ======= ADDITIONAL FONTS IN INRSTEX ==========

% =====special fonts ========
\font\tendu = cmdunh10 
%\font\inchhigh=cminch
%\font\fortysc=cmsc40
%\font\fortypen=pen40
%\font\fortypeni=peni40
%\font\thirtysc=cmsc30
%\font\twentysc=cmsc20
%\font\twentyfourozub=ozub24 %OCR types 
%\font\twentyfourozuh=ozuh24  % ocr hollow
%\font\twentynaro = cmnaro

% ===== twenty one point  =======
\font\twentyonermsca = cmr10 scaled 2074



% ===== eighteen point =======
\font\eighteenrm= cmr17 % no eighteen point font in cm  
\font\eighteenrmsca = cmr10 scaled 1728
%\font\eighteenark = ark18 % hands and pencils
%\font\eighteenbbb = bbb18 % ??????
%\font\eighteenpen = pen18 % pen script like font 
%\font\eighteenpeni = peni18 % pen script like font 

% ===== fourteenpoint fonts =====
\font\fourteenrmsca = cmr10 scaled 1440
\font\fourteenbfxsca = cmbx10 scaled 1440


% ===== twelvepoint fonts 
\font\twelverm= cmr12
\font\twelvermsca = cmr10 scaled 1200
\font\twelvebfxsca = cmbx10 scaled 1200
\let\twelvebf=\twelvebfxsca
\font\twelvess  = cmss12
\font\twelvessb= cmssbx10 scaled 1200
\font\twelveit = cmti12
\font\twelvei = cmmi12
\font\twelvesl = cmsl12
\font\twelvesy = cmsy10 scaled 1200 % no 12point symbol
\font\twelvett = cmtt12
\font\twelvesc = cmcsc10 scaled 1200 % no 12point form 
\let\sanss = \twelvess
\let\sanssb= \tenss
\let\twelveex=\tenex

% ======= tenpoint fonts not in plain.tex ======= 
\font\tenss = cmss10 
\font\tenssb = cmssbx10  % 
\font\tencsc = cmcsc10 % caps and small caps
\font\tenu=cmu10 % unslanted text italic
\font\tenbfx = cmbx10
\font\tenbfy = cmbsy10 %bold math symbols?


% ======== smaller fonts not in plain.tex ========
\font\ninerm=cmr9
\font\eightrm=cmr8
\font\sixrm=cmr6

\font\ninei=cmmi9
\font\eighti=cmmi8
\font\sixi=cmmi6
\skewchar\ninei='177 \skewchar\eighti='177 \skewchar\sixi='177

\font\ninesy=cmsy9
\font\eightsy=cmsy8
\font\sixsy=cmsy6
\skewchar\ninesy='60 \skewchar\eightsy='60 \skewchar\sixsy='60

\font\eightss=cmssq8

\font\eightssi=cmssqi8

\font\ninebfx=cmbx9
\font\eightbfx=cmbx8
\let\sevenbfx = \sevenbf % defined in plain as cmbx7
\font\sixbfx=cmbx6
\let\fivebfx = \fivebf % defined in plain as cmbx5

\font\ninett=cmtt9
\font\eighttt=cmtt8

\hyphenchar\tentt=-1 % inhibit hyphenation in typewriter type
\hyphenchar\twelvett=-1 
\hyphenchar\ninett=-1
\hyphenchar\eighttt=-1

\font\ninesl=cmsl9
\font\eightsl=cmsl8

\font\nineit=cmti9
\font\eightit=cmti8
\font\sevenit= cmti7

% =========== margin note/version  font ===========
\let\notefont = \eighttt

% ============ THE FONT FAMILIES =========== 

% =============  <....>point families =============
% \pointsize and \documentpointsize are the point sizes of the 
% present and document point sizes.
% resets the default spacing 

\newskip\ttglue
\def\twelvepoint{\def\pointsize{twelve}%
  \def\rm{\fam0\twelverm}%
  \textfont0=\twelverm \scriptfont0=\tenrm \scriptscriptfont0=\sevenrm
  \textfont1=\twelvei \scriptfont1=\teni \scriptscriptfont1=\seveni
  \textfont2=\twelvesy \scriptfont2=\tensy \scriptscriptfont2=\sevensy
  \textfont3=\twelveex \scriptfont3=\twelveex \scriptscriptfont3=\twelveex
  \def\it{\fam\itfam\twelveit}%
  \textfont\itfam=\twelveit
  \def\sl{\fam\slfam\twelvesl}%
  \textfont\slfam=\twelvesl
  \def\bf{\fam\bffam\twelvebfxsca}%
  \textfont\bffam=\twelvebfxsca \scriptfont\bffam=\tenbf
   \scriptscriptfont\bffam=\sevenbfx
  \def\tt{\fam\ttfam\twelvett}%
  \textfont\ttfam=\twelvett \scriptfont\ttfam=\tentt
  \tt \ttglue=.5em plus .25em minus .15em
  \normalbaselines\rm % now so resets are ok. 
  \setstrut
  \ifd@ocsty \def\documentpointsize{twelve}%
              \let\footerfont =\twelveit
              \let\headerfont=\twelveit 
              \let\footnotefont = \tenpoint
              \let\smallerfont = \tenpoint
              \let\cheadfont = \twentyonermsca
              \let\sheadfont=\fourteenbfxsca
              \let\ssheadfont=\twelvebfxsca
              \let\sssheadfont=\twelvebfxsca
              \let\dsssheadfont=\twelvebfxsca
              \let\captionnumfont=\twelvebfxsca
              \let\captiontitlefont=\relax
              \let\captionbodyfont=\relax
              \let\eqnumfont=\rm
              \setspacings\fi
  \let\bigfont = \fourteenrmsca
  \let\biggfont = \eighteenrmsca
  \let\bigggfont = \twentyonermsca 
  \let\titlefont = \bigggfont
  \let\sc=\tenrm
  }


% ======= elevenpoint fonts not in plain.tex ======= 

\font\elevenssb = cmssbx10 scaled 1095  % no ss bold
\font\elevencsc = cmcsc10 scaled 1095 % caps and small caps
\font\elevenbfx = cmbx10 scaled 1095 % bold wide
\font\elevenbfy = cmbsy10 scaled 1095 %bold math symbols?

% ===== elevenpoint fonts 
\font\elevenrm= cmr10 scaled 1095
\font\elevenrmsca = cmr10 scaled 1095
\font\elevenbf = cmb10 scaled 1095
\font\elevenbfxsca = cmbx10 scaled 1095
\font\elevenss = cmss10 scaled 1095 
\font\elevenit = cmti10 scaled 1095
\font\eleveni = cmmi10 scaled 1095
\font\elevensl = cmsl10 scaled 1095
\font\elevensy = cmsy10 scaled 1095
\font\eleventt = cmtt10 scaled 1095
\font\elevensc = cmcsc10 scaled 1095
\let\sanss = \elevenss

\let\elevenex=\tenex


\def\elevenpoint{\def\pointsize{eleven}%
  \def\rm{\fam0\elevenrm}%
  \textfont0=\elevenrm \scriptfont0=\eightrm \scriptscriptfont0=\sixrm
  \textfont1=\eleveni \scriptfont1=\eighti \scriptscriptfont1=\sixi
  \textfont2=\elevensy \scriptfont2=\eightsy \scriptscriptfont2=\sixsy
  \textfont3=\elevenex \scriptfont3=\elevenex \scriptscriptfont3=\elevenex
  \def\it{\fam\itfam\elevenit}%
  \textfont\itfam=\elevenit
  \def\sl{\fam\slfam\elevensl}%
  \textfont\slfam=\elevensl
  \def\bf{\fam\bffam\elevenbf}%
  \textfont\bffam=\elevenbf \scriptfont\bffam=\eightbfx
   \scriptscriptfont\bffam=\sixbfx
  \def\tt{\fam\ttfam\eleventt}%
  \textfont\ttfam=\eleventt \scriptfont\ttfam=\tentt
  \tt \ttglue=.5em plus .25em minus .15em
  \normalbaselines\rm % now so resets are ok. 
  \setstrut
  \ifd@ocsty \def\documentpointsize{eleven}%
              \let\footerfont=\elevenit
              \let\headerfont=\elevenit 
              \let\footnotefont = \ninepoint
              \let\smallerfont = \ninepoint
              \let\cheadfont = \eighteenrmsca
              \let\sheadfont=\fourteenbfxsca
              \let\ssheadfont=\elevenbfxsca
              \let\sssheadfont=\elevenbfxsca
              \let\dsssheadfont=\elevenbfxsca
              \let\captionnumfont=\elevenbfxsca
              \let\captiontitlefont=\relax
              \let\captionbodyfont=\relax
              \let\eqnumfont=\rm
              \setspacings\fi
  \let\bigfont = \twelvermsca
  \let\biggfont = \fourteenrmsca
  \let\bigggfont = \eighteenrmsca
  \let\titlefont = \bigggfont
  \let\sc=\ninerm
  }


\def\tenpoint{\def\pointsize{ten}\def\rm{\fam0\tenrm}%
  \textfont0=\tenrm \scriptfont0=\sevenrm \scriptscriptfont0=\fiverm
  \textfont1=\teni \scriptfont1=\seveni \scriptscriptfont1=\fivei
  \textfont2=\tensy \scriptfont2=\sevensy \scriptscriptfont2=\fivesy
  \textfont3=\tenex \scriptfont3=\tenex \scriptscriptfont3=\tenex
  \def\it{\fam\itfam\tenit}%
  \textfont\itfam=\tenit
  \def\sl{\fam\slfam\tensl}%
  \textfont\slfam=\tensl
  \def\bf{\fam\bffam\tenbfx}%
  \textfont\bffam=\tenbfx \scriptfont\bffam=\sevenbfx
   \scriptscriptfont\bffam=\fivebfx
  \def\tt{\fam\ttfam\tentt}%
  \textfont\ttfam=\tentt
  \tt \ttglue=.5em plus .25em minus .15em
  \normalbaselines\rm
  \setstrut
  \ifd@ocsty \def\documentpointsize{ten}\let\footerfont =\tenit 
              \let\headerfont=\tenit 
              \let\footnotefont = \eightpoint 
              \let\smallerfont = \eightpoint
              \let\cheadfont = \eighteenrm 
              \let\sheadfont=\twelvebfxsca
              \let\ssheadfont=\tenbfx
              \let\sssheadfont=\tenbfx
              \let\dsssheadfont=\tenbfx 
              \let\captionnumfont=\tenbfx
              \let\captiontitlefont=\relax
              \let\captionbodyfont=\relax
              \let\eqnumfont=\rm
              \setspacings\fi
  \let\bigfont = \twelvermsca
  \let\biggfont = \fourteenrmsca
  \let\bigggfont = \eighteenrmsca
  \let\titlefont = \bigggfont
  \let\sc=\eightrm }

\def\ninepoint{\def\pointsize{nine}\def\rm{\fam0\ninerm}%
  \textfont0=\ninerm \scriptfont0=\sixrm \scriptscriptfont0=\fiverm
  \textfont1=\ninei \scriptfont1=\sixi \scriptscriptfont1=\fivei
  \textfont2=\ninesy \scriptfont2=\sixsy \scriptscriptfont2=\fivesy
  \textfont3=\tenex \scriptfont3=\tenex \scriptscriptfont3=\tenex
  \def\it{\fam\itfam\nineit}%
  \textfont\itfam=\nineit
  \def\sl{\fam\slfam\ninesl}%
  \textfont\slfam=\ninesl
  \def\bf{\fam\bffam\ninebfx}%
  \textfont\bffam=\ninebfx \scriptfont\bffam=\sixbfx
   \scriptscriptfont\bffam=\fivebfx
  \def\tt{\fam\ttfam\ninett}%
  \textfont\ttfam=\ninett
  \tt \ttglue=.5em plus .25em minus .15em
  \normalbaselines\rm
  \setstrut
  \ifd@ocsty \def\documentpointsize{nine}\let\footerfont =\nineit 
              \let\headerfont=\nineit 
              \let\footnotefont = \eightpoint
              \let\smallerfont = \eightpoint
              \let\cheadfont = \eighteenrm 
              \let\sheadfont=\tenbfx
              \let\ssheadfont=\tenbfx
              \let\sssheadfont=\tenbfx
              \let\dsssheadfont=\tenbfx 
              \let\captionnumfont=\tenbfx
              \let\captiontitlefont=\relax
              \let\captionbodyfont=\relax
              \let\eqnumfont=\rm
              \setspacings\fi
  \let\bigfont = \twelvermsca
  \let\biggfont = \fourteenrmsca
  \let\bigggfont = \eighteenrm 
  \let\titlefont = \bigggfont
  \let\sc=\sevenrm }


\def\eightpoint{\def\pointsize{eight}\def\rm{\fam0\eightrm}%
  \textfont0=\eightrm \scriptfont0=\fiverm \scriptscriptfont0=\fiverm
  \textfont1=\eighti \scriptfont1=\fivei \scriptscriptfont1=\fivei
  \textfont2=\eightsy \scriptfont2=\fivesy \scriptscriptfont2=\fivesy
  \textfont3=\tenex \scriptfont3=\tenex \scriptscriptfont3=\tenex
  \def\it{\fam\itfam\eightit}%
  \textfont\itfam=\eightit
  \def\sl{\fam\slfam\eightsl}%
  \textfont\slfam=\eightsl
  \def\bf{\fam\bffam\eightbfx}%
  \textfont\bffam=\eightbfx \scriptfont\bffam=\fivebfx
   \scriptscriptfont\bffam=\fivebfx
  \def\tt{\fam\ttfam\eighttt}%
  \textfont\ttfam=\eighttt
  \tt \ttglue=.5em plus .25em minus .15em
  \normalbaselines\rm
  \setstrut
  \ifd@ocsty \def\documentpointsize{eight}\let\footerfont =\eightit
                  \let\headerfont=\eightit 
                  \let\footnotefont = \eightpoint 
                  \let\smallerfont = \eightpoint
                  \let\cheadfont = \fourteenrmsca
                  \let\sheadfont=\eightbfx
                  \let\ssheadfont=\eightbfx
                  \let\sssheadfont=\eightbfx
                  \let\dsssheadfont=\eightbfx
                  \let\captionnumfont=\eightbfx
                  \let\captiontitlefont=\relax
                  \let\captionbodyfont=\relax
                  \let\eqnumfont=\rm
                  \setspacings\fi
  \let\bigfont = \tenrm
  \let\biggfont = \twelverm
  \let\bigggfont = \fourteenrmsca
  \let\titlefont = \bigggfont
  \let\sc=\sevenrm }

